\documentclass[12pt,a4paper]{article}

% --- Metadata ---
\def\courseshortheader{HỌC MÁY ỨNG DỤNG --- HỒI QUY LOGISTIC}
\def\coverbookline{TRÍ TUỆ NHÂN TẠO}
\def\covermaintitle{Hồi quy logistic}
\def\coversubtitle{Lý thuyết --- Sigmoid, log loss, chuẩn hóa}

% Shared preamble for nhom_5 (requires \courseshortheader before \input)
\usepackage[utf8]{inputenc}
\usepackage[vietnamese]{babel}
\usepackage{amsmath,amssymb,amsthm}
\usepackage{geometry}
\usepackage{enumitem}
\usepackage[most]{tcolorbox}
\usepackage{hyperref}
\usepackage{fancyhdr}
\usepackage{graphicx}
\usepackage{booktabs}
\usepackage{tikz}
\usetikzlibrary{calc,arrows.meta,decorations.pathreplacing,positioning}
\usepackage{eso-pic}
\usepackage{xcolor}

\geometry{a4paper, margin=2cm, bottom=2.5cm}
\hypersetup{colorlinks=true, linkcolor=blue!70!black, urlcolor=blue}
\setlist[itemize]{topsep=4pt, itemsep=2pt}
\setlist[enumerate]{topsep=4pt, itemsep=2pt}

\AddToShipoutPictureFG{%
  \AtPageCenter{%
    \begin{tikzpicture}[remember picture, overlay]
      \node[rotate=45, scale=1, opacity=0.06]{%
        \fontsize{60}{72}\selectfont\bfseries\sffamily Đặng Tấn Quí%
      };
    \end{tikzpicture}%
  }%
}

\pagestyle{fancy}
\fancyhf{}
\fancyhead[L]{\small\textit{\courseshortheader}}
\fancyhead[R]{\small\textit{Đặng Tấn Quí}}
\fancyfoot[C]{\thepage}
\fancyfoot[L]{\footnotesize\textcopyright\ Đặng Tấn Quí}
\fancyfoot[R]{\footnotesize SĐT: 0377\,122\,210}
\renewcommand{\headrulewidth}{0.4pt}
\renewcommand{\footrulewidth}{0.4pt}

\fancypagestyle{plain}{%
  \fancyhf{}
  \fancyfoot[C]{\thepage}
  \fancyfoot[L]{\footnotesize\textcopyright\ Đặng Tấn Quí}
  \fancyfoot[R]{\footnotesize SĐT: 0377\,122\,210}
  \renewcommand{\headrulewidth}{0pt}
  \renewcommand{\footrulewidth}{0.4pt}
}

\newtcolorbox{dongco}[1][]{%
  enhanced, breakable,
  colback=teal!4!white, colframe=teal!60!black,
  fonttitle=\bfseries, title={Động cơ học -- Tại sao cần khái niệm này?}, #1%
}
\newtcolorbox{ynghia}[1][]{%
  enhanced, breakable,
  colback=purple!3!white, colframe=purple!50!black,
  fonttitle=\bfseries, title={Ý nghĩa \& Trực giác}, #1%
}
\newtcolorbox{ungdung}[1][]{%
  enhanced, breakable,
  colback=olive!5!white, colframe=olive!60!black,
  fonttitle=\bfseries, title={Ứng dụng thực tế}, #1%
}
\newtcolorbox{dinhnghia}[1][]{%
  enhanced, breakable,
  colback=blue!3!white, colframe=blue!60!black,
  fonttitle=\bfseries, title={Định nghĩa}, #1%
}
\newtcolorbox{dinhly}[1][]{%
  enhanced, breakable,
  colback=green!3!white, colframe=green!50!black,
  fonttitle=\bfseries, title={Định lý}, #1%
}
\newtcolorbox{tinhchat}[1][]{%
  enhanced, breakable,
  colback=yellow!5!white, colframe=orange!50!black,
  fonttitle=\bfseries, title={Tính chất}, #1%
}
\newtcolorbox{phuongphap}[1][]{%
  enhanced, breakable,
  colback=gray!5!white, colframe=gray!50!black,
  fonttitle=\bfseries, title={Phương pháp}, #1%
}
\newtcolorbox{luuy}[1][]{%
  enhanced, breakable,
  colback=red!3!white, colframe=red!50!black,
  fonttitle=\bfseries, title={Lưu ý}, #1%
}
\newtcolorbox{congthuc}[1][]{%
  enhanced, breakable,
  colback=cyan!3!white, colframe=cyan!50!black,
  fonttitle=\bfseries, title={Công thức}, #1%
}
\newtcolorbox{vidu}[1][]{%
  enhanced, breakable,
  colback=violet!3!white, colframe=violet!50!black,
  fonttitle=\bfseries, title={Ví dụ}, #1%
}
\newtcolorbox{phanvidu}[1][]{%
  enhanced, breakable,
  colback=red!2!white, colframe=red!40!black,
  fonttitle=\bfseries, title={Phản ví dụ}, #1%
}
\newtcolorbox{tongket}[1][]{%
  enhanced, breakable,
  colback=blue!4!white, colframe=blue!70!black,
  fonttitle=\bfseries, title={Tổng kết chương}, #1%
}


\usepackage{tabularx}
\usepackage{array}
\usepackage{float}
\usepackage{amsmath}

\graphicspath{{images/logistic-regression/}}

\newcommand{\mlimg}[2][0.88]{%
  \def\mlimgfile{images/logistic-regression/#2}%
  \IfFileExists{\mlimgfile}{%
    \begin{center}\includegraphics[width=#1\linewidth]{#2}\end{center}%
  }{%
    \begin{luuy}[title={Thiếu hình minh họa}]%
      Không tìm thấy \expandafter\path\expandafter{\mlimgfile}. Chạy script \path{AI/1_machine_learning/_template/download_mlcc_images.sh}.%
    \end{luuy}%
  }%
}

\begin{document}

\thispagestyle{empty}
\begin{tikzpicture}[remember picture, overlay]
  \fill[blue!80!black] (current page.south west) rectangle (current page.north east);
  \fill[blue!60!cyan, opacity=0.8]
    (current page.south west) -- (current page.north west) --
    ([xshift=-3cm]current page.north east) -- (current page.south east) -- cycle;
  \foreach \r/\o in {6/0.05, 4.5/0.08, 3/0.12} {
    \fill[white, opacity=\o] ([xshift=2cm, yshift=-2cm]current page.north east) circle (\r cm);
  }
  \draw[white, line width=2pt] ([yshift=-5cm]current page.north west) ++(2cm,0) -- ++(17cm,0);
  \draw[white, line width=2pt] ([yshift=-16cm]current page.north west) ++(2cm,0) -- ++(17cm,0);
  \node[anchor=west, white, font=\fontsize{28}{34}\bfseries\selectfont]
    at ([xshift=3cm, yshift=-7.5cm]current page.north west) {\coverbookline};
  \node[anchor=west, white, font=\fontsize{20}{26}\selectfont]
    at ([xshift=3cm, yshift=-9cm]current page.north west) {\covermaintitle};
  \node[anchor=west, white, font=\fontsize{16}{22}\selectfont]
    at ([xshift=3cm, yshift=-10.2cm]current page.north west) {\coversubtitle};
  \node[anchor=west, white, font=\Large\bfseries]
    at ([xshift=3cm, yshift=-13cm]current page.north west) {Biên soạn:};
  \node[anchor=west, yellow!80!white, font=\fontsize{24}{30}\bfseries\selectfont]
    at ([xshift=3cm, yshift=-14.5cm]current page.north west) {ĐẶNG TẤN QUÍ};
  \node[anchor=west, white!90, font=\large]
    at ([xshift=3cm, yshift=-18cm]current page.north west) {SĐT: 0377\,122\,210};
  \node[anchor=south, white!80, font=\small]
    at ([yshift=1.5cm]current page.south) {Tài liệu có bản quyền --- Cấm sao chép khi chưa được sự đồng ý của tác giả};
  \node[anchor=south east, white, font=\Large\bfseries]
    at ([xshift=-2cm, yshift=3cm]current page.south east) {\the\year};
\end{tikzpicture}

\newpage
\tableofcontents
\newpage

\section*{Lời nói đầu}
\addcontentsline{toc}{section}{Lời nói đầu}

\begin{ynghia}[title={Nguồn và cách đọc}]
Tài liệu biên soạn và dịch từ \emph{Logistic Regression} (Google ML Crash Course)\footnote{\url{https://developers.google.com/machine-learning/crash-course/logistic-regression}}.
\end{ynghia}

\begin{luuy}[title={Điều kiện tiên quyết}]
Cần nắm \href{https://developers.google.com/machine-learning/intro-to-ml}{Introduction to Machine Learning} và \href{https://developers.google.com/machine-learning/crash-course/linear-regression}{Linear regression}.
\end{luuy}

\begin{phuongphap}[title={Mục tiêu học tập}]
\begin{itemize}
  \item Nhận diện bài toán phù hợp với hồi quy logistic.
  \item Giải thích cách mô hình dùng sigmoid để tính xác suất.
  \item So sánh hồi quy tuyến tính và hồi quy logistic.
  \item Giải thích vì sao logistic dùng log loss thay vì squared loss.
  \item Giải thích tầm quan trọng của chuẩn hóa khi huấn luyện logistic.
\end{itemize}
\end{phuongphap}

\newpage

%==============================================================================
\section{Giới thiệu hồi quy logistic}
%==============================================================================

\begin{dongco}[title={Từ MPG sang ``có mưa không?''}]
Trong \href{https://developers.google.com/machine-learning/crash-course/linear-regression}{hồi quy tuyến tính}, bạn dự đoán \textbf{số liên tục} (ví dụ MPG). Nhiều câu hỏi thực tế là dạng có/không hoặc \textbf{xác suất}: ``Hôm nay có mưa không?'', ``Email này có phải spam không?''
\end{dongco}

\begin{dinhnghia}[title={Hồi quy logistic}]
\href{https://developers.google.com/machine-learning/glossary\#logistic_regression}{Hồi quy logistic (logistic regression)} dự đoán \textbf{xác suất} của một kết quả (thường nhị phân).
\end{dinhnghia}

\begin{ungdung}[title={Video giới thiệu}]
\url{https://www.youtube.com/watch?v=72AHKztZN44}
\end{ungdung}

\newpage

%==============================================================================
\section{Tính xác suất bằng hàm sigmoid}
%==============================================================================

\begin{dongco}[title={Xác suất 93,2\% hay lớp ``spam''?}]
Nhiều bài toán cần \textbf{xác suất} làm đầu ra. Logistic là cơ chế rất hiệu quả. Bạn có thể:
\begin{enumerate}
  \item Dùng trực tiếp: \texttt{0.932} $\Rightarrow$ 93,2\% khả năng spam.
  \item Chuyển thành lớp: \texttt{Spam} / \texttt{Not Spam} (module \href{https://developers.google.com/machine-learning/crash-course/classification}{Phân loại}).
\end{enumerate}
Phần này dùng xác suất \textbf{trực tiếp}; chuyển sang lớp học ở module sau.
\end{dongco}

\subsection{Hàm sigmoid}

\begin{dinhnghia}[title={Hàm sigmoid}]
Họ \textbf{hàm logistic} cho đầu ra trong $(0,1)$. Hàm chuẩn --- \href{https://developers.google.com/machine-learning/glossary\#sigmoid-function}{sigmoid} (hình chữ S):
\[
  f(x) = \frac{1}{1 + e^{-x}}
\]
\begin{itemize}
  \item $f(x)$: đầu ra sigmoid; $x$: đầu vào.
  \item $e \approx 2{,}71828$ (số Euler).
\end{itemize}
\end{dinhnghia}

\mlimg[0.82]{sigmoid_function_with_axes.png}
\begin{center}\small\textbf{Hình 1.} Đồ thị sigmoid: tiến về 0 khi $x \to -\infty$, tiến về 1 khi $x \to +\infty$ (không đạt đúng 0 hay 1).\end{center}

\begin{vidu}[title={Bảng giá trị sigmoid ($x$ từ $-7$ đến $7$)}]
\begin{center}
\small
\begin{tabular}{cc|cc}
\textbf{$x$} & \textbf{$f(x)$} & \textbf{$x$} & \textbf{$f(x)$} \\
\midrule
$-7$ & 0.001 & 0 & 0.500 \\
$-6$ & 0.002 & 1 & 0.731 \\
$-5$ & 0.007 & 2 & 0.881 \\
$-4$ & 0.018 & 3 & 0.952 \\
$-3$ & 0.047 & 4 & 0.982 \\
$-2$ & 0.119 & 5 & 0.993 \\
$-1$ & 0.269 & 6 & 0.997 \\
 & & 7 & 0.999 \\
\end{tabular}
\end{center}
Dù $|x|$ lớn đến đâu, đầu ra luôn trong $(0,1)$.
\end{vidu}

\subsection{Biến đổi đầu ra tuyến tính thành xác suất}

\begin{congthuc}[title={Thành phần tuyến tính và dự đoán logistic}]
\[
  z = b + w_1 x_1 + w_2 x_2 + \cdots + w_N x_N,\qquad
  y' = \frac{1}{1 + e^{-z}}
\]
\begin{itemize}
  \item $z$: \href{https://developers.google.com/machine-learning/glossary\#log-odds}{\textbf{log odds}} (logit).
  \item $b$: bias; $w_i$: trọng số học được; $x_i$: đặc trưng của ví dụ.
  \item $y'$: xác suất dự đoán trong $(0,1)$.
\end{itemize}
\end{congthuc}

\begin{ynghia}[title={Log-odds}]
Từ $y = \dfrac{1}{1 + e^{-z}}$ suy ra $z = \ln\!\left(\dfrac{y}{1-y}\right)$ --- logarit tỷ số xác suất hai kết quả.
\end{ynghia}

\mlimg[0.92]{linear_to_logistic.png}
\begin{center}\small\textbf{Hình 2.} Trái: $z = 2x + 5$. Phải: cùng ba điểm sau sigmoid; ví dụ $z=-10 \Rightarrow y' \approx 0{,}00004$ (ô vàng).\end{center}

\begin{tinhchat}[title={Tuyến tính vs sigmoid}]
Phương trình tuyến tính có thể cho $z$ rất lớn/nhỏ; sigmoid ``uốn'' đầu ra về $(0,1)$, không bao gồm hai đầu.
\end{tinhchat}

\begin{luuy}[title={Kiểm tra hiểu biết --- Tính $z$ và $y'$}]
Cho $b=1$, $w_1=2$, $w_2=-1$, $w_3=5$; $x_1=0$, $x_2=10$, $x_3=2$.

\medskip
\textbf{Câu 1 --- $z$?} \textbf{Đúng: $z=1$.} Công thức $z = 1 + 2x_1 - x_2 + 5x_3$ $\Rightarrow$ $1+0-10+10=1$. Sai: $0{,}731$ là $y'$, không phải $z$.

\medskip
\textbf{Câu 2 --- $y'$?} \textbf{Đúng: $\approx 0{,}731$.}
\[
  y' = \frac{1}{1+e^{-1}} \approx \frac{1}{1{,}367} \approx 0{,}731
\]
\end{luuy}

\newpage

%==============================================================================
\section{Mất mát và chuẩn hóa}
%==============================================================================

\begin{dongco}[title={Cùng quy trình huấn luyện, khác loss và chuẩn hóa}]
Logistic huấn luyện giống tuyến tính (gradient descent), nhưng:
\begin{enumerate}
  \item Dùng \href{https://developers.google.com/machine-learning/glossary\#Log_Loss}{\textbf{Log Loss}}, không dùng \href{https://developers.google.com/machine-learning/glossary\#l2-loss}{squared loss}.
  \item \href{https://developers.google.com/machine-learning/glossary\#regularization}{Chuẩn hóa} gần như bắt buộc để tránh \href{https://developers.google.com/machine-learning/glossary\#overfitting}{overfitting}.
\end{enumerate}
\end{dongco}

\subsection{Log Loss}

\begin{ynghia}[title={Vì sao không dùng squared loss?}]
Squared loss phù hợp khi tốc độ thay đổi $y'$ gần hằng (ví dụ $y'=b+3x_1$). Với sigmoid (chữ S), gần $z=0$ thì $y'$ đổi mạnh; khi $|z|$ lớn thì đổi ít. Gần 0/1, squared loss cần lưu sai số rất nhỏ $\Rightarrow$ tốn bộ nhớ. Log loss dùng \textbf{logarit độ lớn thay đổi}.
\end{ynghia}

\begin{vidu}[title={Độ chính xác cần thiết khi $z$ lớn}]
\begin{center}
\small
\begin{tabular}{ccc}
\textbf{Đầu vào} & \textbf{Đầu ra logistic} & \textbf{Số chữ số có nghĩa} \\
\midrule
5 & 0.993 & 3 \\
7 & 0.999 & 3 \\
9 & 0.9999 & 4 \\
10 & 0.99998 & 5 \\
\end{tabular}
\end{center}
\end{vidu}

\begin{congthuc}[title={Log Loss (trung bình trên $N$ mẫu)}]
\[
  \text{Log Loss} = -\frac{1}{N}\sum_{i=1}^{N}\Bigl[y_i\log(y_i') + (1-y_i)\log(1-y_i')\Bigr]
\]
\begin{itemize}
  \item $N$: số mẫu có nhãn; $y_i \in \{0,1\}$; $y_i' \in (0,1)$.
  \item Dùng \textbf{trung bình} (không phải tổng) để tách biệt tinh chỉnh batch size và learning rate.
\end{itemize}
\end{congthuc}

\subsection{Chuẩn hóa trong hồi quy logistic}

\begin{dinhnghia}[title={Regularization}]
\href{https://developers.google.com/machine-learning/glossary\#regularization}{Chuẩn hóa} phạt độ phức tạp khi huấn luyện. Không chuẩn hóa, logistic có thể cứ đẩy loss về 0 khi có nhiều đặc trưng. Thường dùng:
\begin{itemize}
  \item \href{https://developers.google.com/machine-learning/crash-course/overfitting/regularization}{Chuẩn hóa $L_2$}
  \item \href{https://developers.google.com/machine-learning/crash-course/overfitting/regularization\#early_stopping_an_alternative_to_complexity-based_regularization}{Early stopping} --- giới hạn số bước huấn luyện.
\end{itemize}
\end{dinhnghia}

\begin{luuy}
Chi tiết thêm: module \href{https://developers.google.com/machine-learning/crash-course/overfitting}{Datasets, Generalization, and Overfitting}.
\end{luuy}

\begin{tongket}[title={Tổng kết chương --- Hồi quy logistic}]
\begin{itemize}
  \item $z = b + \sum w_i x_i$; $y' = \sigma(z) \in (0,1)$.
  \item Huấn luyện: Log Loss + chuẩn hóa ($L_2$ hoặc early stopping).
  \item Tiếp theo: \href{https://developers.google.com/machine-learning/crash-course/classification}{Phân loại} --- ngưỡng chuyển xác suất thành lớp.
\end{itemize}
\end{tongket}

\end{document}
